\documentclass[
  spanish,
  aspectratio=169,
  xcolor={dvipsnames,table}
]{beamer}

\geometry{paperwidth=19.2cm,paperheight=10.8cm}
\usepackage[utf8]{inputenc}
\usepackage[T1]{fontenc}
\usepackage[spanish,es-tabla]{babel}
\usepackage[scaled=.96]{helvet}
\renewcommand{\familydefault}{\sfdefault}
\usepackage{booktabs}
\usepackage{graphicx}
\usepackage{ragged2e}
\usepackage{tikz}
\usepackage{hyperref}
\usepackage{listings}
\usetikzlibrary{calc,positioning,fit,shadows.blur}

% Estilo de bloque de código para reproducir la participación textual del foro
\lstdefinestyle{foroplano}{
  basicstyle=\ttfamily\scriptsize\color{unadmInk},
  backgroundcolor=\color{unadmPaper},
  breaklines=true,
  columns=fullflexible,
  keepspaces=true,
  extendedchars=true,
  inputencoding=utf8,
  frame=leftline,
  framerule=1.4pt,
  rulecolor=\color{unadmGreen},
  xleftmargin=0.5em,
  aboveskip=0.3em,
  belowskip=0.3em,
  literate=
    {á}{{\'a}}1 {é}{{\'e}}1 {í}{{\'i}}1 {ó}{{\'o}}1 {ú}{{\'u}}1
    {Á}{{\'A}}1 {É}{{\'E}}1 {Í}{{\'I}}1 {Ó}{{\'O}}1 {Ú}{{\'U}}1
    {ñ}{{\~n}}1 {Ñ}{{\~N}}1 {¿}{{?`}}1 {¡}{{!`}}1
}

\newcommand{\studentname}{Martin Jonathan de la Cruz}
\newcommand{\studentshort}{M. J. de la Cruz}
\newcommand{\studentid}{ES2611202040}
\newcommand{\universityname}{Universidad Abierta y a Distancia de México}
\newcommand{\facultyname}{Licenciatura en Derecho}
\newcommand{\coursename}{Garantías Constitucionales}
\newcommand{\coursecode}{LDE-S2B1}
\newcommand{\activitytitle}{Foro diagnóstico}
\newcommand{\activitysubtitle}{Nociones básicas de derecho constitucional}
\newcommand{\documentsubject}{Actividad 1 - Garantías Constitucionales}
\newcommand{\activityweek}{Semana 1}
\newcommand{\teachingfigure}{Aarón López Pérez}
\newcommand{\deliverydate}{\today}
\newcommand{\locationname}{Roma Norte, Ciudad de México}
\newcommand{\departmentlogo}{img/departamentos/UnADM.pdf}

\definecolor{unadmGreenDark}{HTML}{174A3A}
\definecolor{unadmGreen}{HTML}{5F8F3A}
\definecolor{unadmGold}{HTML}{B88A2A}
\definecolor{unadmPaper}{HTML}{F6F7F2}
\definecolor{unadmInk}{HTML}{1F2A24}

\mode<presentation>{
  \usetheme{default}
  \usefonttheme{professionalfonts}
  \setbeamertemplate{navigation symbols}{}
  \setbeamertemplate{blocks}[rounded][shadow=false]
}
\setbeamersize{text margin left=0.60cm,text margin right=0.60cm}
\setbeamercolor{background canvas}{bg=white}
\setbeamercolor{normal text}{fg=unadmInk,bg=white}
\setbeamercolor{structure}{fg=unadmGreenDark}
\setbeamercolor{frametitle}{fg=white,bg=unadmGreenDark}
\setbeamercolor{block title}{fg=white,bg=unadmGreen}
\setbeamercolor{block body}{fg=unadmInk,bg=unadmPaper}
\setbeamerfont{title}{size=\LARGE,series=\bfseries}
\setbeamerfont{frametitle}{size=\large,series=\bfseries}
\setbeamertemplate{itemize item}{\textcolor{unadmGreen}{\large$\blacktriangleright$}}

\setbeamertemplate{frametitle}{
  \nointerlineskip
  \begin{beamercolorbox}[wd=\textwidth,ht=1.02cm,dp=0.22cm,leftskip=0.35cm,rightskip=0.25cm]{frametitle}
    \insertframetitle\hfill\includegraphics[height=0.60cm]{\departmentlogo}
  \end{beamercolorbox}
  {\color{unadmGold}\rule{\textwidth}{1.3pt}}
}

\setbeamertemplate{footline}{
  \leavevmode
  \hbox{
    \begin{beamercolorbox}[wd=.34\paperwidth,ht=2.7ex,dp=1ex,leftskip=1em]{author in head/foot}\color{white}\studentshort\end{beamercolorbox}
    \begin{beamercolorbox}[wd=.40\paperwidth,ht=2.7ex,dp=1ex,center]{title in head/foot}\color{white}\coursename\end{beamercolorbox}
    \begin{beamercolorbox}[wd=.26\paperwidth,ht=2.7ex,dp=1ex,rightskip=1em plus 1fil]{date in head/foot}\color{white}\coursecode\hfill\insertframenumber/\inserttotalframenumber\end{beamercolorbox}
  }
}
\setbeamercolor{author in head/foot}{bg=unadmGreenDark}
\setbeamercolor{title in head/foot}{bg=unadmGreen}
\setbeamercolor{date in head/foot}{bg=unadmGreenDark}

\title[\coursecode]{\activitytitle}
\subtitle{\activitysubtitle}
\author[\studentshort]{\studentname\\\footnotesize Matrícula: \studentid}
\institute[UnADM]{\universityname\\\facultyname}
\date[\activityweek]{\locationname\\\activityweek\\\deliverydate}

\begin{document}

\begin{frame}[plain]
  \begin{tikzpicture}[remember picture,overlay]
    \fill[unadmGreenDark] (current page.south west) rectangle ([xshift=0.58\paperwidth]current page.north west);
    \fill[unadmGreen] ([xshift=0.58\paperwidth]current page.south west) rectangle (current page.north east);
    \node[anchor=center,opacity=0.10] at (current page.center) {\includegraphics[width=0.72\paperwidth]{\departmentlogo}};
    \node[anchor=north west] at ([xshift=0.58cm,yshift=-0.46cm]current page.north west) {\includegraphics[width=2.55cm]{\departmentlogo}};
  \end{tikzpicture}
  \vspace*{1.62cm}
  {\color{white}\fontsize{24}{29}\selectfont\bfseries \activitytitle\par}
  \vspace{0.28cm}
  {\color{white!86}\large \activitysubtitle}\\[0.16cm]
  {\color{white!82}\normalsize \documentsubject}\\[0.35cm]
  {\color{unadmGold}\rule{0.58\linewidth}{1.3pt}}\\[0.35cm]
  {\color{white}\studentname}\\
  {\color{white!82}\footnotesize \facultyname}
\end{frame}

\begin{frame}{Propósito del foro}
  \begin{block}{Diagnóstico inicial}
    Identificar conocimientos previos, experiencias, expectativas y necesidades de aprendizaje sobre nociones jurídicas básicas.
  \end{block}
  \begin{itemize}
    \item No sustituye la enseñanza: orienta la planeación.
    \item Permite detectar confusiones conceptuales tempranas.
    \item Vincula participación activa con nivelación jurídica.
  \end{itemize}
\end{frame}

\begin{frame}{Constitución}
  \begin{block}{Respuesta propia}
    Es la norma fundamental que organiza el poder público, reconoce derechos y limita la actuación de las autoridades.
  \end{block}
  \begin{itemize}
    \item Da validez al resto del sistema jurídico.
    \item Protege derechos frente a actos contrarios a la Constitución.
    \item Es el punto de partida de las garantías constitucionales.
  \end{itemize}
\end{frame}

\begin{frame}{Derecho objetivo y subjetivo}
  \begin{columns}[T]
    \begin{column}{0.48\linewidth}
      \begin{block}{Derecho objetivo}
        Conjunto de normas jurídicas vigentes.
      \end{block}
    \end{column}
    \begin{column}{0.48\linewidth}
      \begin{block}{Derecho subjetivo}
        Facultad o pretensión concreta de una persona.
      \end{block}
    \end{column}
  \end{columns}
  \vspace{0.25cm}
  La regla pertenece al plano objetivo; la posibilidad de exigir algo pertenece al plano subjetivo.
\end{frame}

\begin{frame}{Derecho sustantivo y adjetivo}
  \begin{itemize}
    \item \textbf{Sustantivo}: define derechos, obligaciones, prohibiciones e instituciones de fondo.
    \item \textbf{Adjetivo o procesal}: establece las vías y formas para hacer valer esos derechos.
  \end{itemize}
  \begin{block}{Diferencia central}
    El primero responde qué derecho existe; el segundo responde cómo se reclama o protege.
  \end{block}
\end{frame}

\begin{frame}{Proceso y procedimiento}
  \begin{block}{Proceso}
    Estructura jurídica completa orientada a resolver una controversia ante autoridad competente.
  \end{block}
  \begin{block}{Procedimiento}
    Secuencia concreta de actos, etapas, plazos y formalidades dentro del proceso.
  \end{block}
\end{frame}

\begin{frame}{Derecho humano y garantía}
  \begin{columns}[T]
    \begin{column}{0.48\linewidth}
      \begin{block}{Derecho humano}
        Prerrogativa vinculada con la dignidad de la persona.
      \end{block}
    \end{column}
    \begin{column}{0.48\linewidth}
      \begin{block}{Garantía}
        Mecanismo jurídico que protege o hace exigible ese derecho.
      \end{block}
    \end{column}
  \end{columns}
  \vspace{0.25cm}
  Sin garantías eficaces, los derechos corren el riesgo de quedar solo como declaraciones formales.
\end{frame}

\begin{frame}{Microcaso de aplicación}
  \begin{block}{Detención arbitraria}
    Si una persona es detenida sin información, sin defensa y sin presentación oportuna ante autoridad, pueden afectarse libertad personal, seguridad jurídica y debido proceso.
  \end{block}
  \begin{itemize}
    \item El derecho sustantivo identifica derechos afectados.
    \item El derecho adjetivo indica la vía para reclamarlos.
    \item La garantía funciona como mecanismo de protección.
  \end{itemize}
\end{frame}

% META (interno): evidencia del producto FORO (100 Técnicas Didácticas, UnADM).
% Reproduce de forma textual la participación publicada en el foro.
\begin{frame}[fragile]{Participación publicada en el foro}
  \begin{lstlisting}[style=foroplano]
Asunto: Foro diagnóstico - Nociones básicas de derecho

Hola a todas y todos: comparto mi participación.

1) Constitución: norma fundamental; organiza el poder, reconoce derechos y limita a la autoridad (arts. 1.o y 133).
2) Objetivo (sistema de normas) vs. subjetivo (facultad concreta que se exige a partir de la norma).
3) Sustantivo (define el fondo) vs. adjetivo (vías para reclamarlo).
4) Proceso (estructura para resolver el conflicto) vs. procedimiento (secuencia de actos, plazos y formalidades).
5) Derecho humano (prerrogativa ligada a la dignidad) vs. garantía (mecanismo que lo hace exigible, p. ej. el juicio de amparo).

Pregunta para el grupo: ¿en qué caso concreto un derecho reconocido se quedó sin protección por falta de una garantía eficaz?
  \end{lstlisting}
\end{frame}

\begin{frame}{Cierre}
  \begin{block}{Aprendizaje principal}
    El foro diagnóstico permite iniciar la asignatura con evidencia sobre qué conceptos domina o confunde el grupo.
  \end{block}
  \begin{itemize}
    \item Ayuda a planificar actividades de nivelación.
    \item Fortalece el pensamiento jurídico inicial.
    \item Prepara el estudio de garantías constitucionales y derechos humanos.
  \end{itemize}
\end{frame}

\end{document}
